\label{section:BoundaryConditions}
\index{Boundary conditions}
Boundary conditions are assigned after the \gwf, \swf\ and \cln\ domains have been generated and cells (and, where needed, nodes) have been selected.
In all cases you must first make a cell selection as described on page~\pageref{page:cellSelect} for the active domain.
This section groups the available instructions by boundary-condition type; under each type the domain-specific forms (\gwf, \swf\ and \cln) are described where they apply.

\subsection{Specified Head (Constant Head)}
\index{Boundary conditions ! constant head}
A constant head (CHD) boundary condition fixes the hydraulic head at a cell, allowing water to flow into or out of the domain depending on surrounding conditions.

\subsubsection{\gwf\ Domain}
Named CHD groups can be reported as separate volumetric-budget terms in the Modflow \texttt{.lst} file (and therefore as separate Tecplot VolumeBudget columns). Immediately before assigning constant head, tag the currently chosen cells with:

\ins{chd zone name}
    {
        \squish
        \begin{enumerate}
        \item \str{ZoneName}   Budget label for this CHD group (maximum 16 characters; underscores recommended).
        \end{enumerate}
          The label is applied to cells that receive the next {\sf gwf constant head} assignment. \mfus\ prints separate IN/OUT lines for each named zone, and still prints the aggregate \texttt{CONSTANT HEAD} total.
    }

To assign a uniform constant head to the \gwf\ model domain use this instruction:

\ins{gwf constant head}
    {
        \squish
        \begin{enumerate}
        \item \rnum{CHead} [$L$]   Constant hydraulic head.
        \end{enumerate}
          A constant hydraulic head of \rnum{CHead} is assigned to the chosen cells.
          The CHD package record writes \rnum{CHead} as both the start and end head. \cln\ and \swf\ constant-head records instead ramp from the cell starting head to \rnum{CHead}; see those instructions and Appendix~\ref{appendix:release_verification}.
    }

\subsubsection{\swf\ Domain}
\ins{swf constant head}
    {
        \squish
        \begin{enumerate}
        \item \rnum{CHead} [$L$]  Constant hydraulic head.
        \end{enumerate}
          A constant hydraulic head of \rnum{CHead} is assigned to the chosen cells.
          The CHD package record writes two heads: the cell starting head at the beginning of the stress period and \rnum{CHead} at the end. \mfus\ interpolates between them so the assigned head is not applied as an instantaneous shock.
    }

To prescribe a head that changes from one stress period to the next (for example a tidal sea level), place this instruction after each {\sf stress period} block:

\ins{swf transient constant head}
    {
        \squish
        \begin{enumerate}
        \item \rnum{StartHead}, \rnum{EndHead} [$L$]  Head at the start and end of the current stress period (one line).
        \end{enumerate}
          The chosen cells receive a CHD record for the current stress period only, with start head \rnum{StartHead} and end head \rnum{EndHead}. \mfus\ interpolates linearly between them over the stress period, so a sequence of short stress periods gives a piecewise-linear head history.
          Stress periods without a {\sf swf transient constant head} assignment reuse the previous CHD list. A preceding {\sf chd zone name} applies to these records. A cell cannot carry both a static {\sf swf constant head} and a transient one.
    }

\subsubsection{\cln\ Domain}
\ins{cln constant head}
    {
        \squish
        \begin{enumerate}
        \item \rnum{CHead} [$L$]  Constant hydraulic head.
        \end{enumerate}
          A constant hydraulic head of \rnum{CHead} is assigned to the chosen cells.
          The CHD package record writes two heads: the cell starting head at the beginning of the stress period and \rnum{CHead} at the end. \mfus\ interpolates between them so the assigned head is not applied as an instantaneous shock.
    }

\subsection{Drains}
\index{Boundary conditions ! drains}
A drain boundary condition allows water to leave the domain when the cell head is above the drain elevation. Currently only the \gwf\ domain supports drains in \mut.

\subsubsection{\gwf\ Domain}
\ins{gwf drain}
    {
        \squish
        \begin{enumerate}
        \item \rnum{DrainK}  [$L$ $T^{-1}$]  Drain conductance.
        \end{enumerate}
          A drain conductance of \rnum{DrainK} is assigned to the chosen cells.  The top elevation of the cell is assigned automatically as the drain elevation.
    }

\subsection{Specified Flux (Wells)}
\index{Boundary conditions ! wells}
A well boundary condition forces a volumetric flux into or out of chosen cells. Positive rates add water; negative rates remove water.

\subsubsection{\gwf\ Domain}
\ins{gwf well}
    {
        \squish
        \begin{enumerate}
            \item \rnum{PumpRate} [$L^{3}$ $T^{-1}$]  Pumping rate.
        \end{enumerate}
        A pumping rate of \rnum{PumpRate} is assigned to the chosen cells. Positive pumping rates add water to the domain, negative pumping rates remove water from the domain.
    }

\subsubsection{\swf\ Domain}
\ins{swf well}
    {
        \squish
        \begin{enumerate}
            \item \rnum{PumpRate} [$L^{3}$ $T^{-1}$]  Pumping rate.
        \end{enumerate}
        A pumping rate of \rnum{PumpRate} is assigned to the chosen cells.  Positive pumping rates add water to the domain, negative pumping rates remove water from the domain.
    }

\subsubsection{\cln\ Domain}
\ins{cln well}
    {
        \squish
        \begin{enumerate}
            \item \rnum{PumpRate} [$L^{3}$ $T^{-1}$]  Pumping rate.
        \end{enumerate}
        A pumping rate of \rnum{PumpRate} is assigned to the chosen cells. Positive pumping rates add water to the domain, negative pumping rates remove water from the domain.
    }

\subsection{Recharge}
\index{Boundary conditions ! recharge}
Recharge applies a flux per unit area to chosen cells.
\mut\ supports constant-rate stress-period \textsf{RCH}, multi-zone transient \textsf{RTS} on the \swf\ domain, and General Spatio-Temporal Recharge (\textsf{GSTR}) on \gwf, \swf\ and \cln.

\subsubsection{Constant-rate recharge (\textsf{RCH})}
\paragraph*{\gwf\ Domain}
\ins{gwf recharge}
    {
        \squish
        \begin{enumerate}
            \item \rnum{RechRate} [$L$ $T^{-1}$]  Recharge rate.
            \item \inum{RechOpt}  Recharge option.
        \end{enumerate}
        A recharge rate of \rnum{RechRate} is assigned to the chosen cells.

        The recharge option \inum{RechOpt} is used to define where the recharge is to be applied and can have one of the following values:
        \begin{enumerate}
            \item To top layer
            \item To one specified node in each vertical column
            \item To highest active node in each vertical column
            \item To the swf domain on top of each vertical column
        \end{enumerate}
    }

\paragraph*{\swf\ Domain}
\ins{swf recharge}
    {
        \squish
        \begin{enumerate}
            \item \rnum{RechRate} [$L$ $T^{-1}$] Recharge rate.
            \item \inum{RechOpt}  Recharge option.
        \end{enumerate}
        A recharge rate of \rnum{RechRate} is assigned to the chosen cells.

        The recharge option \inum{RechOpt} is used to define where the recharge is to be applied and in this case should be set to a value of 4, which applies the recharge to the \swf\ domain.
    }

Constant-rate \textsf{RCH} is written per stress period.
For time-varying rates on the \swf\ domain, use {\sf swf transient recharge} below.
Spatially and temporally varying ArcASCII recharge that coexists with \textsf{RCH}/\textsf{RTS} is documented under General Spatio-Temporal Recharge (Section~\ref{section:GSTR}).

\subsection{Evapotranspiration (\textsf{EVT})}
\index{Boundary conditions ! evapotranspiration}
\index{Boundary conditions ! EVT}
Classical MODFLOW \textsf{EVT} removes water from the \gwf\ saturated zone as a function of water-table depth. The maximum rate applies when head is at or above the ET surface; the rate declines linearly to zero over the extinction depth below that surface. \textsf{EVT} does not evaporate ponded water from the \swf\ domain.

\paragraph*{\gwf\ Domain}
\ins{gwf evt}
    {
        \squish
        \begin{enumerate}
            \item \rnum{EVTR} [$L$ $T^{-1}$]  Maximum ET rate.
            \item \inum{NEVTOP}  ET option (1 = top layer; 3 = highest active cell in each column). Option 2 is not yet supported by \mut.
            \item \rnum{EXDP} [$L$]  Extinction depth below the ET surface (must be $> 0$).
        \end{enumerate}
        The ET surface elevation (\texttt{SURF}) is set automatically to the land-surface elevation (top of layer~1) for each vertical column.
        Chosen cells receive \rnum{EVTR}; unchosen top-layer cells receive a zero rate.
        Repeat the instruction after each \textsf{end stress period} to change rates through the simulation.
        \mut\ writes \texttt{Modflow.evt} and registers \textsf{EVT} in the name file.
    }

\subsubsection{Transient recharge (\textsf{RTS})}
\index{Boundary conditions ! transient recharge}
Transient multi-zone recharge is currently supported for the \swf\ domain.

\paragraph*{\swf\ Domain}
For time-varying recharge rates supplied by an RTS file, use this instruction instead of a constant {\sf swf recharge}:

\ins{swf transient recharge}
    {
        \squish
        \begin{enumerate}
            \item \str{FNameRTS} Name of a single-column RTS file (rates in model length/time units).
            \item \inum{RechOpt} Recharge option (use 4 for the \swf\ domain).
        \end{enumerate}
        Currently chosen cells are tagged with the next RTS zone index.
        Repeat the instruction with a different cell selection and RTS file to define additional zones.
        Each source file must contain exactly one rate column; \mut\ merges the zones into a multi-column file \texttt{\_buildo.merged\_RTS.rts} and writes an RTS-capable RCH package with matching \texttt{INRCHZONES} / \texttt{IZNRCH} maps.
        Record counts and time stamps must align across zone files.
        Build and post-processing \texttt{.eco} files report whether recharge used stress-period \textsf{RCH} or an RTS file; post-processing also writes applied-rate Tecplot output (e.g.\ \texttt{*.SWF\_RCH\_RTS.tecplot.dat}).
    }

\subsubsection{General Spatio-Temporal Recharge (\textsf{GSTR})}
\label{section:GSTR}
The \mfus\ \textsf{GSTR} package assigns recharge that varies in space and time using ArcASCII raster snapshots.
It operates at the timestep level, independently of the existing \textsf{RCH} and \textsf{RTS} packages.
Existing \textsf{RCH}/\textsf{RTS} input files remain fully supported; when both \textsf{GSTR} and \textsf{RCH} (or \textsf{RTS}) are active, recharge fluxes from each package are added to the flow equation.

\paragraph*{MUT instructions}
Select the cells that belong to the instance, optionally name the instance, then assign snapshots on the target domain.

\ins{gstr instance name}
    {
    \squish
    \begin{enumerate}
    \item \str{name} Budget label for the instance (at most 16 characters).
    \end{enumerate}
    The name is applied to the next \texttt{gwf gstr}, \texttt{swf gstr}, or \texttt{cln gstr} instruction.
    If omitted, MUT assigns a default name \texttt{GSTR\_$n$}.
}

\ins{gwf gstr}
{
    \squish
    \begin{enumerate}
    \item \rnum{FAC} Multiplier on sampled rates (use $1$ for none).
    \item One or more lines \rnum{time} \str{raster.asc} (ascending times, model time units).
    \item \texttt{end gstr}
    \end{enumerate}
    Currently chosen \gwf\ cells are written with global node number and centroid $(x,y)$.
    Identical forms exist for \swf\ (\texttt{swf gstr}, IDOMAIN$=3$) and \cln\ (\texttt{cln gstr}, IDOMAIN$=2$).
}

Example (SWF rainfall with start and end rasters):
\begin{verbatim}
active domain
swf
    choose all cells
    gstr instance name
    Rain_SWF
    swf gstr
    1.0
    0.0    rain_t0.asc
    20.0   rain_t20.asc
    end gstr
\end{verbatim}

MUT writes \texttt{\textit{prefix}.gstr} and a name-file entry \texttt{GSTR}.
A KAERI/KURT demonstration lives under
\fpath{KURT_Model/2_models/2_TestTransientRainfall/4_GSTR}.

\paragraph*{Name file}
Add one name-file entry:
\begin{verbatim}
GSTR  <unit>  <prefix>.gstr
\end{verbatim}

\paragraph*{Input file format}
Comment lines may begin with \texttt{\#} or \texttt{!}. Blank lines are skipped.

\begin{verbatim}
NGSTR  IGSTRCB
INSTANCE  <name>
  IDOMAIN  NCELLS  NSNAPS  [FAC]
  <inode>  <x>  <y>
  ...
  <time>   <raster.asc>
  ...
INSTANCE  <name2>
  ...
\end{verbatim}

\begin{description}
\item[\texttt{NGSTR}] Number of named GSTR instances.
\item[\texttt{IGSTRCB}] Unit number for cell-by-cell budget terms (\texttt{0} = none).
\item[\texttt{INSTANCE} \textit{name}] Starts a named instance. The name (up to 16 characters) is used as the volumetric budget label for that instance.
\item[\texttt{IDOMAIN}] Domain code: \texttt{1}=\gwf, \texttt{2}=\cln, \texttt{3}=\swf. Used for validation; recharge is applied at the listed global node numbers.
\item[\texttt{NCELLS}] Number of cells in the instance (all domain cells or a subset).
\item[\texttt{NSNAPS}] Number of time snapshots (must be $\ge 1$).
\item[\texttt{FAC}] Optional multiplier applied to sampled recharge rates (default $1$).
\item[\texttt{inode x y}] Global node number and cell coordinates used to sample rasters. Coordinates must be in the same length units and coordinate system as the ArcASCII grids.
\item[\texttt{time raster.asc}] Snapshot time (model time units, ascending) and path to an ArcASCII recharge raster ($L/T$).
\end{description}

\paragraph*{ArcASCII rasters}
Each snapshot file is a standard ArcASCII grid with header keywords
\texttt{ncols}, \texttt{nrows}, \texttt{xllcorner} or \texttt{xllcenter},
\texttt{yllcorner} or \texttt{yllcenter}, \texttt{cellsize}, and optional \texttt{nodata\_value},
followed by row-major cell values (north to south).
Bilinear interpolation is used at each cell $(x,y)$.
Samples outside the raster extent, or involving a nodata corner, are treated as zero recharge.

\paragraph*{Time interpolation}
At each timestep, \mfus\ uses elapsed simulation time (\texttt{TOTIM}) to select the raster snapshots immediately before and after the current time and linearly interpolates the sampled rate for each cell.
Outside the snapshot range, the nearest endpoint snapshot is held (no extrapolation beyond the first or last raster).
The volumetric flux applied to a node is
\[
Q = \mathrm{FAC}\times r(x,y,t)\times A_n
\]
where $r$ is the interpolated rate ($L/T$) and $A_n$ is the node area stored by \mfus\ (GWF area, SWF face area, or CLN conduit area as appropriate).

\paragraph*{Budget}
Each named instance contributes its own budget term using the instance name as the label.
\textsf{GSTR} does not replace or modify \textsf{RCH}/\textsf{RTS} budget terms.

\paragraph*{Post-processing to Tecplot}
After \texttt{postprocess existing modflow model}, MUT reads the \textsf{GSTR} package and reconstructs applied rates by sampling the snapshot rasters (and a mid-interval interpolated zone when two or more snaps exist). Outputs include:
\begin{verbatim}
_posto.o.<prefix>.SWF_GSTR.tecplot.dat
_posto.o.<prefix>.SWF_GSTR_TS.tecplot.dat
\end{verbatim}
(and analogous \texttt{GWF\_GSTR} / \texttt{CLN\_GSTR} files by domain).
Named instance totals also appear in the listing-file volume-budget Tecplot export.


\subsection{Critical Depth Outflow}
\index{Boundary conditions ! critical depth}
A critical depth boundary condition assigned to a \swf\ cell allows water to flow out of the domain at a rate that depends on the surface water depth and a contributing length (the length of the cell side over which the outflow occurs).
This boundary type applies only to the \swf\ domain.

\subsubsection{\swf\ Domain}
One of the following two instructions can be used:

\ins{swf critical depth with sidelength1}
    {
          A critical depth outflow boundary condition is assigned to the chosen cells.

          It is assumed that an accurate estimate of the contributing length of a cell can be based on the square root of the cell's horizontal area.
    }

\ins{swf critical depth}
    {
          A critical depth outflow boundary condition is assigned to the chosen cells.

          The contributing length of the cell outflow boundary is calculated from the \swf\ mesh outer boundary nodes, with each outer boundary node connected to a chosen cell contributing a half-element side length in both directions along the outer boundary.
    }

Although {\sf swf critical depth} is less convenient than {\sf swf critical depth with sidelength1}, it does calculate a contributing length that matches the actual length along the \swf\ mesh outer boundary.

\paragraph{Named outlets.}
\index{Boundary conditions ! critical depth ! named zones}
By default all critical depth cells are reported together as a single \texttt{SWBC} volumetric-budget term. To report separate outlets (for example two catchment outlets), tag each group of cells immediately before the critical depth instruction with:

\ins{swbc zone name}
    {
        \squish
        \begin{enumerate}
        \item \str{ZoneName}   Budget label for this critical depth group (maximum 16 characters, no blanks; underscores recommended).
        \end{enumerate}
          The label is applied to the cells that receive the next {\sf swf critical depth} or {\sf swf critical depth with sidelength1} assignment. Repeating a name adds cells to the existing zone.
    }

When at least one zone is named, \mut\ writes a zone id as a third column of the \texttt{.swbc} file and a \texttt{.swbczone} map file (NAM type \texttt{SWBZ}). \mfus\ then prints one IN/OUT line per zone instead of \texttt{SWBC}, and saves one \swf\ cell-by-cell record per zone; critical depth cells without a zone keep the \texttt{SWBC} label. The sum of the zone terms equals the unzoned \texttt{SWBC} term. Post-processing adds \texttt{IN\_}/\texttt{OUT\_}{\em ZoneName} columns to the VolumeBudget Tecplot file, \swf\ results variables \swf\ to {\em ZoneName} alongside the combined \swf\ to SWBC, and a \texttt{SWBC Zone} column to the build scatter file \texttt{*\_SWBC.tecplot.dat}.

For example, two outlets A and B:
\begin{verbatim}
        clear chosen cells
        choose cells from xyz list
            224562.618075775,329692.764796124,0
            224524.557045581,329512.523755688,0
            end
        swbc zone name
        Outlet_A
        swf critical depth

        clear chosen cells
        choose cells from xyz list
            235922.820107068,330050.46088884,0
            235799.630126383,330094.549842411,0
            end
        swbc zone name
        Outlet_B
        swf critical depth
\end{verbatim}

The example \texttt{6\_Abdul\_Prism\_Cell} uses the second approach:
\begin{verbatim}
        clear chosen nodes
        choose gb nodes
        ./gb/grid.nchos.Outer boundary nodes
        flag chosen nodes as outer boundary

        clear chosen cells
        clear chosen nodes
        choose cells from gb elements
        ./gb/grid.echos.Critical depth outlet
        swf critical depth
\end{verbatim}

Some key features of this example are:
\begin{itemize}
    \item \swf\ {\em nodes} are chosen and flagged to be on the outer boundary with the instruction {\sf flag chosen nodes as outer boundary} (see node selection on page~\pageref{page:nodeSelect}).
    \item \swf\ {\em cells} are chosen using a \gb\ chosen {\em elements} file with the instruction {\sf choose cells from gb elements}.  Because this example was generated using the mesh-centred control volume approach, there is a 1-to-1 correspondence between template mesh elements and \swf\ cells.

        In the example \texttt{6\_Abdul\_Prism\_Cell\_nc}, the node-centred control volume approach is used and the instruction {\sf choose cells from gb nodes} is used instead, because in this case there is a 1-to-1 correspondence between template mesh {\em nodes} and \swf\ cells.
\end{itemize}
